\paraid{p-a26667c35d82}
The Gromov--Hausdorff distance
$\GHDistance$
quantifies the difference in shape between
two compact metric spaces.
Let
$\GHSpace$
denote the space of isometry classes of nonempty compact metric spaces,
equipped with the Gromov--Hausdorff distance
$\GHDistance$.
We call
the space
$(\GHSpace, \GHDistance)$
the \emph{Gromov--Hausdorff space}.
Antonyan
\cite[p.~2]{Antonyan2020Euclidean}
asked whether
$\GHSpace$
 is homeomorphic to
$\SeparableHilbertSpace$.
We answer this question affirmatively by proving the following classification.

\begin{theorem}[\Cref{thm:part-iv-main}]\label{thm:main}
\paraid{p-f4cb7088c681}
The space
$\GHSpace$
is homeomorphic to the Hilbert space
$\SeparableHilbertSpace$.
\end{theorem}

\paraid{p-414311e7f43d}
Zava
\cite[arXiv v3, Theorem~B]{Zava2025Coarse}
proves that the subspace of
$\GHSpace$
consisting of finite metric spaces cannot be coarsely embedded into any
Hilbert space.
In particular,
no homeomorphism in \Cref{thm:main} can be bi-Lipschitz.

\medskip
\noindent\textbf{Related work.}
\paraid{p-fb73239bb9f0}
For comparison with \Cref{thm:main},
we first recall a classification of subspaces of
$\GHSpace$.
For each positive integer
$\FiniteSize $,
Antonyan identifies the subspace of
$\GHSpace$
represented by compact subsets of
$\RealNumbers^\FiniteSize $
with the Hilbert cube minus a point
\cite[Corollary~5.3]{Antonyan2021Euclidean}.

\paraid{p-31fbde650eca}
Geodesics and optimal correspondences describe the metric structure
underlying our topological classification of
$\GHSpace$.
Ivanov,
Nikolaeva,
and Tuzhilin
\cite[arXiv v1, Theorem~1]{IvanovNikolaevaTuzhilin2015}
prove that any two points of
$\GHSpace$
can be joined by a geodesic.
Ivanov,
Iliadis,
and Tuzhilin
\cite[arXiv v1, Theorem~2.6 and Corollary~2.7]{IvanovIliadisTuzhilin2016}
prove the existence of optimal correspondences between compact metric spaces
and use them to attain the Gromov--Hausdorff distance as the Hausdorff distance
between isometric copies in a common metric space.
A further point of comparison for our classification by homeomorphisms
is rigidity under surjective isometries.
Ivanov and Tuzhilin
\cite[Main Theorem]{IvanovTuzhilin2019Isometries}
prove that every surjective isometry of
$\GHSpace$
is the identity.
Tuzhilin surveys Hausdorff and Gromov--Hausdorff distance geometry in
\cite{Tuzhilin2020}.

\paraid{p-9f8f911356f3}
Earlier work studies the topology of
$\GHSpace$
through continuous families and embedded compact spaces.
Our result identifies the topological type of
$\GHSpace$
itself.
The author constructed families of branching Gromov--Hausdorff geodesics
continuously parametrized by the Hilbert cube
\cite[Theorem~1.3]{Ishiki2022Branching}.
The author also constructed topological embeddings of arbitrary compact metrizable
spaces into the subspaces of
$\GHSpace$
consisting of continua
\cite[Theorem~1.1]{Ishiki2022Continua},
spaces with prescribed dimensions
\cite[Theorem~1.3]{Ishiki2023FractalDimensions},
and compact metric trees
\cite[Theorem~1.1]{Ishiki2023MetricTrees}.
In each case,
the embeddings can be chosen with finitely many distinct prescribed values.
Byakuno
\cite[arXiv v1, Corollary~1.4]{Byakuno2026Embedding}
proves that countable products of closed intervals with positive summable lengths,
equipped with the supremum metric,
admit isometric embeddings into
$\GHSpace$.

\paraid{p-6f8cb7cfa910}
For comparison with our classification of the full space
$\GHSpace$,
we recall local descriptions and dimension formulas for its subspaces
of finite metric spaces.
Iliadis,
Ivanov,
and Tuzhilin
\cite[arXiv v1, Theorem~4.1]{IliadisIvanovTuzhilin2017Local}
prove that sufficiently small neighborhoods of generic finite metric spaces,
within the subspace of spaces with the same number of points,
are isometric to open subsets of a finite-dimensional space with the maximum norm.
Here generic means that the nonzero distances are pairwise distinct
and all triangle inequalities involving three distinct points are strict.
Nakajima,
Yamauchi,
and Zava
\cite[arXiv v1, Theorem~A]{NakajimaYamauchiZava2025}
prove that the subspace of spaces with at most
$\FiniteSize$
points has topological dimension
$\FiniteSize(\FiniteSize-1)/2$.
A complementary approach to the global topology of
$\GHSpace$
studies its compactifications.
Nakajima and Shioya
\cite[arXiv v1, Main Theorem~1.2]{NakajimaShioyaCompactification}
construct a compact metrizable space containing a dense topological copy of
$\GHSpace$.

\paraid{p-d8629ae38ad9}
Part~\ref{part:measures} selects probability measures continuously over
$\GHSpace$.
Topological results for metric measure spaces provide a comparison when
the measure is included as part of the data.
Kazukawa,
Nakajima,
and Shioya
\cite[arXiv v1, Theorems~1.4--1.6]{KazukawaNakajimaShioya2024}
prove that the space of metric measure spaces is contractible and locally
path connected both in the box topology and in the concentration topology.
Our classification concerns compact metric spaces.
An extension of the Gromov--Hausdorff framework treats noncomplete spaces
together with their boundary.
Shibahara
\cite[arXiv v1, Definition~3.2 and Theorem~1.1]{Shibahara2021Boundary}
introduces a Gromov--Hausdorff metric with boundary on the isometry classes
of noncomplete precompact locally compact metric spaces.
The boundary is the complement of the space in its metric completion.
Our local models in Part~\ref{part:spectral} keep track of isometries
between representatives.
Isometries between individual spaces are also retained by the moduli stack
used by
Yuji
\cite[arXiv v2, Definition~5.1]{Yuji2026ModuliStack}
to study families of compact metric spaces.


\paraid{p-86969a2d1c01}
Part~\ref{part:topology} studies
$\GHSpace$
as a global analogue of a hyperspace with a fixed ambient metric space.
The link between these two settings also appears in the study of geodesics.
M\'emoli and Wan
\cite[arXiv v2, Theorem~1]{MemoliWan2023}
prove that every Gromov--Hausdorff geodesic can be realized as a Hausdorff
geodesic of compact subsets of one compact metric space.
For comparison with our absolute retract theorem,
we recall the case of compact subsets of a fixed Banach space.
Curtis
\cite[Theorem~1.6]{Curtis1980}
proves,
in particular,
that the hyperspace of nonempty compact subsets of a Banach space,
equipped with the Hausdorff metric of its norm,
is an absolute retract.
We construct approximations to the identity of
$\GHSpace$
using quotients of hyperspaces of nonempty compact subsets of Banach spaces
by compact groups.

\paraid{p-5f89de56c669}
Our use of Toru\'nczyk's characterization in Part~\ref{part:hilbert}
has a precedent in the study of the Urysohn universal metric space.
In the early twenty-first century,
Uspenskij
\cite[arXiv v1, Theorem~2.1]{Uspenskij2004Urysohn}
proved that this space is homeomorphic to
$\SeparableHilbertSpace$.
The Urysohn space is the complete separable metric space that contains
an isometric copy of every separable metric space and in which every isometry
between finite subsets extends to a surjective isometry of the whole space.


\medskip
\noindent\textbf{Proof strategy and the four parts.}
\paraid{p-ce4b777e5ddb}
We first prove that
$\GHSpace$
is an absolute retract for metrizable spaces.
We approximate a sequence of maps from compact metrizable domains so that
the images of the approximating maps form a discrete family.
Since
$\GHSpace$
is complete and separable,
the absolute retract property and the discrete approximation property
imply \Cref{thm:main} by Toru\'nczyk's characterization
\cite[p.~248, assertion~(i)]{Torunczyk1981},
using the correction in
\cite[Section~C]{Torunczyk1985}.
In Parts~\ref{part:measures}--\ref{part:topology},
we construct the approximations needed for the absolute retract property.
The discrete approximation argument in Part~\ref{part:hilbert}
uses only the common preliminaries and its own metric constructions.
Figure~\ref{fig:proof-strategy}
shows how these two arguments combine to prove the classification.

\paraid{p-69b86b1ef543}
A difficulty in the absolute retract argument is that the isometry group
depends on the compact metric space.
Rouyer's results on generic compact metric spaces
\cite[arXiv v1, Theorems~2 and~4]{Rouyer2011}
imply that
\[
 \overline{
 \{\MetricSpace{\BaseCarrier}{\MetricSymbol}\in\GHSpace
   \mid\Isom\MetricSpace{\BaseCarrier}{\MetricSymbol}
      =\{\IdentityMap_\BaseCarrier\}\}
 }=\GHSpace.
\]
This also implies that the assignment of isometry groups
with their uniform metrics is discontinuous as a map from
$\GHSpace$
to itself.
Our construction must also handle spaces with nontrivial isometry groups.
On each model neighborhood,
we construct homomorphisms from the varying isometry groups into one fixed
orthogonal group so that the local point maps are equivariant.
We then use a partition of unity to combine the local point maps in a Banach space.
A countable product of orthogonal groups acts on this Banach space.
Passing to the orbit space of nonempty compact subsets of this Banach space
makes the resulting approximation independent of all coordinate choices.

\medskip
\noindent\textbf{Part~\ref{part:measures}.
Continuous invariant probability measures.}
\paraid{p-b0e31260869e}
In the first part,
we construct measures for integration on varying compact spaces.
The main result of this part,
\Cref{thm:invariant-measures},
assigns a Borel probability measure
$\SelectedMeasure{\BaseCarrier}{\MetricSymbol}$
to every nonempty compact metric space
$\MetricSpace{\BaseCarrier}{\MetricSymbol}$.
Each assigned measure has full support,
\[
 \supp(\SelectedMeasure{\BaseCarrier}{\MetricSymbol})=\BaseCarrier.
\]
For every isometry
$\IdentifyingIsometry\colon
\MetricSpace{\BaseCarrier}{\MetricSymbol}\to
\MetricSpace{\ComparisonCarrier}{\ComparisonMetric}$,
the assignment satisfies
\[
 \IdentifyingIsometry\Pushforward\SelectedMeasure{\BaseCarrier}{\MetricSymbol}
 =\SelectedMeasure{\ComparisonCarrier}{\ComparisonMetric}.
\]
The assignment of measures is also continuous in the following sense.
If compact subspaces
$\BaseCarrier_\SequenceIndex$
and
$\BaseCarrier$
of a compact metric space
$\MetricSpace{\AmbientCarrier}{\AmbientMetric}$
carry the restricted metrics
$\MetricSymbol_\SequenceIndex$
and
$\MetricSymbol$,
then
\[
 \HausdorffDistance{\AmbientMetric}(\BaseCarrier_\SequenceIndex,\BaseCarrier)\to0
 \quad\Longrightarrow\quad
 \SelectedMeasure{\BaseCarrier_\SequenceIndex}{\MetricSymbol_\SequenceIndex}
 \to\SelectedMeasure{\BaseCarrier}{\MetricSymbol}
 \quad\text{weakly in }\ProbabilityMeasures(\AmbientCarrier).
\]
Here the measures are regarded as probabilities on the common ambient space.

\paraid{p-968906a0f31d}
For a fixed compact space,
averaging over its isometry group produces an invariant probability measure.
Our construction must also preserve continuity when the space varies.
To construct such measures,
we prove that the space of invariant full-support measured compact spaces
is Polish and that the map that forgets the measure is an open surjection
onto
$\GHSpace$.
By Valov's theorem (\Cref{thm:valov}),
we obtain continuous probability laws on its fibers.
Averaging these laws produces the selected measures.
We use full support and continuity in the spectral approximation of
Part~\ref{part:spectral}.

\medskip
\noindent\textbf{Part~\ref{part:spectral}.
Continuous finite-dimensional approximation.}
\paraid{p-16a5923a2a64}
In the second part,
we replace distance functions by finite-dimensional data while controlling
all pairwise distances.
The main result of this part is \Cref{thm:local-models}.
For every
$\MetricSpace{\BaseCarrier}{\MetricSymbol}\in\GHSpace$
and every
$\LocalTolerance>0$,
we obtain a neighborhood
$\ModelNeighborhood$
of
$\MetricSpace{\BaseCarrier}{\MetricSymbol}$
and an integer
$\ModelDimension\geq1$
with the following property.
For each compact metric space
$\MetricSpace{\ComparisonCarrier}{\ComparisonMetric}$
whose isometry class belongs to
$\ModelNeighborhood$,
we construct a norm
$\NormSymbol_\ComparisonCarrier$
on
$\RealNumbers^\ModelDimension$
and a continuous map
$\FeatureCoordinates_\ComparisonCarrier\colon
\ComparisonCarrier\to\RealNumbers^\ModelDimension$.
For every
$\ComparisonPoint,\IntegrationPoint\in\ComparisonCarrier$,
define the induced pseudometric by
\[
 \Pseudometric_\ComparisonCarrier(\ComparisonPoint,\IntegrationPoint)
 =
 \NormSymbol_\ComparisonCarrier
 \bigl(\FeatureCoordinates_\ComparisonCarrier(\ComparisonPoint)
       -\FeatureCoordinates_\ComparisonCarrier(\IntegrationPoint)\bigr).
\]
After shrinking
$\ModelNeighborhood$,
we obtain
\[
 \sup_{\ComparisonPoint,\IntegrationPoint\in\ComparisonCarrier}
 \bigl|\Pseudometric_\ComparisonCarrier(\ComparisonPoint,\IntegrationPoint)
       -\ComparisonMetric(\ComparisonPoint,\IntegrationPoint)\bigr|
 <\LocalTolerance.
\]
The dimension
$\ModelDimension$
is fixed on
$\ModelNeighborhood$,
but may depend on the center
$\MetricSpace{\BaseCarrier}{\MetricSymbol}$
and the error scale
$\LocalTolerance$.
The norms
$\NormSymbol_\ComparisonCarrier$
and point maps
$\FeatureCoordinates_\ComparisonCarrier\colon
\ComparisonCarrier\to\RealNumbers^\ModelDimension$
vary continuously up to changes of coordinates in
$\OrthogonalGroup(\ModelDimension)$
in the sense of
\Cref{thm:local-models}\ref{item:model-continuity}.
Each isometry group
$\Isom\MetricSpace{\ComparisonCarrier}{\ComparisonMetric}$
acts on the coordinate image
$\FeatureCoordinates_\ComparisonCarrier(\ComparisonCarrier)
\subset\RealNumbers^\ModelDimension$
through a continuous homomorphism
\[
 \IsometryRepresentation{\ComparisonCarrier}{\ComparisonMetric}
 \colon\Isom\MetricSpace{\ComparisonCarrier}{\ComparisonMetric}
 \to\OrthogonalGroup(\ModelDimension).
\]
For every
$\FirstIsometry\in\Isom\MetricSpace{\ComparisonCarrier}{\ComparisonMetric}$,
it satisfies
\[
 \FeatureCoordinates_\ComparisonCarrier\circ\FirstIsometry
 =\IsometryRepresentation{\ComparisonCarrier}{\ComparisonMetric}(\FirstIsometry)
  \circ\FeatureCoordinates_\ComparisonCarrier,
 \qquad
 \NormSymbol_\ComparisonCarrier
  \circ\IsometryRepresentation{\ComparisonCarrier}{\ComparisonMetric}(\FirstIsometry)
 =\NormSymbol_\ComparisonCarrier
\]
by Remark~\ref{rem:local-isometry-representation}.
Thus the induced pseudometrics
$\Pseudometric_\ComparisonCarrier$
respect isometries while the coordinate changes take place in the fixed
compact group
$\OrthogonalGroup(\ModelDimension)$.

\paraid{p-458636bfa314}
To construct $\ModelNeighborhood$, we use the measures of Part~\ref{part:measures}
to define the distance operator
$\DistanceOperator{\BaseCarrier}{\MetricSymbol}$
on
$\LebesgueSpace^2(\BaseCarrier,\SelectedMeasure{\BaseCarrier}{\MetricSymbol})$.
For every
$f\in\LebesgueSpace^2(\BaseCarrier,\SelectedMeasure{\BaseCarrier}{\MetricSymbol})$
and
$\BasePoint\in\BaseCarrier$,
set
\[
 (\DistanceOperator{\BaseCarrier}{\MetricSymbol}f)(\BasePoint)
 =\int_\BaseCarrier
   \MetricSymbol(\BasePoint,\IntegrationPoint)f(\IntegrationPoint)
   \,\IntegrationDifferential
   \SelectedMeasure{\BaseCarrier}{\MetricSymbol}(\IntegrationPoint).
\]
These operators and their spectral coordinates were studied by Maria,
Oudot,
and Solomon
\cite[Section~3]{MariaOudotSolomon2020}.
Finite spectral subspaces of the distance operators approximate the distance
functions.
For our local models,
we choose the unique best approximations in
$\LebesgueSpace^\IntegrabilityExponent$
for a finite exponent
$2\leq\IntegrabilityExponent<\infty$
chosen to achieve the error bound at the center in
\ref{item:model-center-error} of \Cref{thm:local-models}.
We keep whole eigenspaces to handle eigenvalue multiplicities by orthogonal
changes of basis.
We retain both the compact coordinate image and its norm for the
construction in Part~\ref{part:topology}.
We expect this approach to provide a bridge between the study of the
Gromov--Hausdorff space and geometric analysis or spectral geometry.

\paraid{p-24f2b5c3820e}
By \cite[Theorem~1.1]{Ishiki2026Interpolation},
we obtain a uniform error bound when extending prescribed metrics from a
discrete family of closed subsets of a fixed metrizable space.
The absolute retract argument here uses approximations of metrics on varying
compact spaces whose coordinate pairs satisfy the convergence conditions in
\ref{item:model-continuity} of \Cref{thm:local-models}.

\medskip
\noindent\textbf{Part~\ref{part:topology}.
The absolute retract property.}
\paraid{p-159ea0252542}
In \Cref{thm:part-iii-main},
the main result of the third part,
we prove that
$\GHSpace$
is an absolute retract for all metrizable spaces.
In \Cref{thm:absolute-extensor},
we prove that
$\GHSpace$
is an absolute extensor for all metrizable spaces.

\paraid{p-02e057f0dca4}
Fix a continuous function
$\ErrorControl\colon\GHSpace\to(0,\infty)$.
We choose a countable locally finite cover by local models whose errors
are less than
$\ErrorControl$
throughout their domains.
For each model index
$\CoverIndex\in\NonnegativeIntegers$,
let
$\ModelDimension_\CoverIndex$
be the model dimension.
We represent the varying norms isometrically in fixed Banach spaces
and combine the point maps by a partition of unity in their
$\SummableNorm$-sum,
a separable Banach space
$\BanachAmbientSpace$.
The norm representation respects orthogonal changes of coordinates
(\Cref{lem:dual-representation,lem:dual-representation-equivariance}).
The compact group
\[
 \ActingGroup
 =\prod_{\CoverIndex\in\NonnegativeIntegers}
   \OrthogonalGroup(\ModelDimension_\CoverIndex)
\]
acts continuously on
$\BanachAmbientSpace$
by linear isometries.
Passing to the orbit of each compact coordinate image removes the
choices of bases and isometric representatives.
The orbit space
$\AuxiliaryAR=\CompactHyperspace(\BanachAmbientSpace)/\ActingGroup$
is an absolute retract by Antonyan's results
(\Cref{lem:hyperspace-orbit-ar}).
By \Cref{thm:global-domination},
we obtain continuous maps
\[
 \GHSpace\xrightarrow{\ \ApproximationMap\ }
 \AuxiliaryAR\xrightarrow{\ \RealizationMap\ }\GHSpace
\]
and a continuous homotopy
$\ApproximationHomotopy\colon\GHSpace\times\UnitInterval\to\GHSpace$.
For every
$\MetricSpace{\ComparisonCarrier}{\ComparisonMetric}\in\GHSpace$,
this homotopy satisfies
\begin{equation}\label{eq:intro-domination-1}
 \ApproximationHomotopy(\MetricSpace{\ComparisonCarrier}{\ComparisonMetric},0)
 =\MetricSpace{\ComparisonCarrier}{\ComparisonMetric},
\end{equation}
and
\begin{equation}\label{eq:intro-domination-2}
 \ApproximationHomotopy(\MetricSpace{\ComparisonCarrier}{\ComparisonMetric},1)
 =\RealizationMap\ApproximationMap\MetricSpace{\ComparisonCarrier}{\ComparisonMetric}.
\end{equation}
The homotopy tracks satisfy
\begin{equation}\label{eq:intro-domination-3}
 \diam_{\GHDistance}
 \bigl\{\ApproximationHomotopy(\MetricSpace{\ComparisonCarrier}{\ComparisonMetric},t)
       \mid t\in\UnitInterval\bigr\}
 <\ErrorControl\MetricSpace{\ComparisonCarrier}{\ComparisonMetric}.
\end{equation}
Applying Hanner's domination theorem (\Cref{thm:hanner-domination})
to the small homotopy tracks in \eqref{eq:intro-domination-3},
we obtain the ANR property.
Scaling all distances to
$0$
contracts
$\GHSpace$
to the one-point space.
Using the ANR/ANE equivalence in \Cref{thm:retract-extensor},
we apply \Cref{thm:contractible-ane} to conclude that
$\GHSpace$
is an AE for metrizable spaces.
By the AE/AR equivalence in \Cref{thm:retract-extensor},
we obtain the absolute retract property.

\begin{figure}[H]
\centering
\begin{tikzpicture}[
 >=Stealth,
 font=\small,
 partbox/.style={
  draw=black,
  fill=white,
  rounded corners=3pt,
  line width=.65pt,
  text width=4.9cm,
  minimum height=2.35cm,
  align=center,
  inner sep=7pt
 },
 route/.style={->,line width=.85pt,draw=black}
]
\node[partbox] (measures) at (-2.95,0) {
 \textbf{Part~\ref{part:measures}}\\
 Continuous invariant\\
 full-support probabilities\\
 \Cref{thm:invariant-measures}
};
\node[partbox] (models) at (2.95,0) {
 \textbf{Part~\ref{part:spectral}}\\
 Finite-dimensional local models\\
 $\Isom(Y,e)\to\OrthogonalGroup(n)$\\
 \Cref{thm:local-models}
};
\node[partbox,minimum height=2.95cm] (ar) at (2.95,-3.55) {
 \textbf{Part~\ref{part:topology}}\\
 $\ActingGroup=\prod_{i\in\NonnegativeIntegers}\OrthogonalGroup(n_i)$\\
 $\CompactHyperspace(\BanachAmbientSpace)/\ActingGroup$ is an AR\\
 Small tracks $\Longrightarrow$ ANR\\
 Scaling contraction $\Longrightarrow$ AR\\
 \Cref{thm:part-iii-main}
};
\node[
 partbox,
 minimum height=2.95cm
] (discrete) at (-2.95,-3.55) {
 \textbf{Part~\ref{part:hilbert}}\\
 Products with finite equilateral spaces\\
 Discrete approximation\\
 \Cref{thm:discrete-approximation}
};
\node[
 partbox,
 text width=10.8cm,
 minimum height=1.9cm
] (hilbert) at (0,-7) {
 \textbf{Toru\'nczyk's characterization}\\
 AR and discrete approximation,
 with completeness and separability\\
 $\GHSpace$ is homeomorphic to real $\SeparableHilbertSpace$\\
 \textbf{Part~\ref{part:hilbert}.
 \Cref{thm:part-iv-main}}
};
\draw[route] (measures.east)--(models.west);
\draw[route] (models.south)--(ar.north)
 node[midway,left=5pt,font=\footnotesize,align=right]
 {Partition of unity\\in the Banach space $\BanachAmbientSpace$};
\draw[route] (ar.south)--(2.95,-5.5)--(0,-5.5)--(hilbert.north);
\draw[line width=.85pt,draw=black]
 (discrete.south)--(-2.95,-5.5)--(0,-5.5);
\fill[black] (0,-5.5) circle (1.5pt);
\end{tikzpicture}
\caption{The two arguments used to identify the topology of
$\GHSpace$.
Arrows indicate dependence of the proofs.
In Parts~\ref{part:measures}--\ref{part:topology},
we prove the absolute retract property.
The discrete approximation theorem in Part~\ref{part:hilbert}
uses the common preliminaries and its own metric constructions.
These two conclusions,
together with completeness and separability,
imply the Hilbert space classification.}
\label{fig:proof-strategy}
\end{figure}


\medskip
\noindent\textbf{Part~\ref{part:hilbert}.
Discrete approximation and Hilbert space topology.}
\paraid{p-64cc2f159292}
In the fourth part,
we prove the remaining approximation condition and then
obtain the classification in \Cref{thm:part-iv-main}.
\Cref{thm:discrete-approximation},
the discrete approximation theorem of this part,
applies to any sequence
$\{\ApproximationDomain_\FamilyIndex\}_{\FamilyIndex\in\NonnegativeIntegers}$
of compact metrizable spaces.
If
$\InputMap_\FamilyIndex\colon\ApproximationDomain_\FamilyIndex\to\GHSpace$
are continuous and
$\OpenCover$
is an open cover of
$\GHSpace$,
then there are continuous maps
$\ApproximatingMap_\FamilyIndex\colon\ApproximationDomain_\FamilyIndex\to\GHSpace$
such that
\[
 \forall\FamilyIndex\in\NonnegativeIntegers\
 \forall x\in\ApproximationDomain_\FamilyIndex\
 \exists U\in\OpenCover,\qquad
 \{\InputMap_\FamilyIndex(x),\ApproximatingMap_\FamilyIndex(x)\}\subset U,
\]
and
\[
 \{\ApproximatingMap_\FamilyIndex(\ApproximationDomain_\FamilyIndex)
   \}_{\FamilyIndex\in\NonnegativeIntegers}
 \text{ is discrete in }\GHSpace.
\]
The last condition means that every point of
$\GHSpace$
has a neighborhood that intersects at most one of these image sets.

\paraid{p-66ca83693c4d}
Taking products with sufficiently small finite equilateral spaces changes each
input by a controlled amount in Gromov--Hausdorff distance.
We choose the cardinalities recursively to exceed the bounds on separated subsets
in the previously constructed compact images.
This separates the images,
and letting the cardinalities tend to infinity prevents accumulation in
$\GHSpace$.
This construction is independent of the results of Parts~\ref{part:measures}--\ref{part:topology}.
If all domains are the Hilbert cube
$\HilbertCube=\UnitInterval^{\NonnegativeIntegers}$,
then we obtain the discrete approximation condition in Toru\'nczyk's theorem.
We combine this condition with completeness,
separability,
and the absolute retract property from Part~\ref{part:topology}
to prove \Cref{thm:main}.

\paraid{p-59ed21a6dca9}
The main theorem also answers the local contractibility question.
We use local contractibility in the sense that every neighborhood
$\OpenNeighborhood $
of a point contains a neighborhood
$\SmallerNeighborhood $
whose inclusion into
$\OpenNeighborhood $
is null-homotopic.
By \Cref{thm:main},
every point has a neighborhood basis consisting of open sets that
strongly deformation retract onto that point.
This conclusion is topological and does not prescribe the neighborhoods
as balls for
$\GHDistance$.

\medskip
\noindent\textbf{Organization.}
\paraid{p-10cf165b0a9f}
The introduction and
\Cref{sec:preliminaries}
provide the common background for four parts.
Each part also has its own preliminaries,
in Sections~\ref{sec:prelim-measures},
\ref{sec:prelim-spectral},
\ref{sec:prelim-topology},
and~\ref{sec:prelim-hilbert},
respectively.
Part~\ref{part:measures},
introduced in \Cref{sec:intro-measures} and developed in \Cref{sec:measures},
constructs continuous invariant probabilities.
Part~\ref{part:spectral},
introduced in \Cref{sec:intro-spectral} and developed in \Cref{sec:spectral},
uses these probabilities to construct finite-dimensional local models.
Part~\ref{part:topology},
introduced in \Cref{sec:intro-topology} and developed in \Cref{sec:ar},
combines the models into global approximations
and proves the absolute retract property.
Part~\ref{part:hilbert},
introduced in \Cref{sec:intro-hilbert} and developed in \Cref{sec:approximation},
proves the discrete approximation property and completes the proof of
\Cref{thm:main}.
\Cref{sec:questions} concludes with questions about related spaces
and the geometry of GH balls.

\medskip
\noindent\textbf{Conventions and notation.}
\paraid{p-67846154e305}
For a set
$\BaseCarrier$,
we denote by
$\Card(\BaseCarrier)$
its cardinality.
Every metric space
$\MetricSpace{\BaseCarrier }{\MetricSymbol }$
representing a point of
$\GHSpace$
is nonempty and compact.
We identify such a pair with its isometry class when it occurs in
$\GHSpace$.
The symbol
$\BaseCarrier $
denotes the underlying set,
equipped with the topology induced by
$\MetricSymbol $.
Subscripts on constructions such as
$\SelectedMeasure{\BaseCarrier}{\MetricSymbol}$
refer to the specified pair
$\MetricSpace{\BaseCarrier }{\MetricSymbol }$.
When isometric embeddings into
$\MetricSpace{\AmbientCarrier }{\AmbientMetric }$
are specified,
we use the same symbols for the embedded subsets and the restricted
metrics.
The notation
$\IdentityMap_{\BaseCarrier}$
denotes the identity map on
$\BaseCarrier$.
We omit the subscript when the domain is specified.
Banach and Hilbert spaces are real unless complex scalars are explicitly indicated.
We denote by
$\CompactHyperspace(\BaseCarrier )$
the set of all nonempty compact subsets of
$\BaseCarrier $.
For a specified metric
$\MetricSymbol $,
we equip it with the Hausdorff metric
$\HausdorffDistance{\MetricSymbol }$.
For a metrizable space
$\DomainSpace$,
we denote by
$\ProbabilityMeasures(\DomainSpace )$
the set of Radon probability measures on
$\DomainSpace$
and equip it with the topology of weak convergence against bounded continuous
functions.
An isometric embedding
$\IsometricEmbeddingMap \colon\MetricSpace{\BaseCarrier }{\MetricSymbol }\to\MetricSpace{\ComparisonCarrier }{\ComparisonMetric }$
is a map satisfying
$\ComparisonMetric (\IsometricEmbeddingMap (\BasePoint ),\IsometricEmbeddingMap (\ComparisonPoint ))=\MetricSymbol (\BasePoint ,\ComparisonPoint )$
for every
$\BasePoint ,\ComparisonPoint \in \BaseCarrier $.
Throughout this paper,
an isometry is a surjective isometric embedding.
We use
$\NonnegativeIntegers=\{0,1,2,\ldots\}$
as the index set of sequences unless stated otherwise.

\begin{useofai}
\paraid{p-64859c0d8f24}
OpenAI Codex,
based on GPT-6,
was used to explore and develop the mathematical constructions and proofs,
search the literature,
and prepare the exposition and \LaTeX{} source.
The research and writing workflow used customized versions of
\texttt{math-research-harness}
(\url{https://github.com/haruhisa-enomoto/math-research-harness})
and
\texttt{math-paper-skills}
(\url{https://github.com/haruhisa-enomoto/math-paper-skills}),
both made publicly available by Haruhisa Enomoto.
The author takes full responsibility for the mathematical content and the final
manuscript.
\end{useofai}

\begin{acknowledgements}
\paraid{p-e60be87ab1f9}
The author was supported by JSPS KAKENHI Grant Number JP24KJ0182.
\end{acknowledgements}
